% 05_evaluation.tex - Seccion de evaluacion experimental.
\section{Experimental Evaluation}

\subsection{Experimental Setup}
% TODO: condiciones, datasets o participantes, metricas y procedimiento.
% Todo valor numerico debe provenir de data/stats_summary.json.

\subsection{Results}
% Tabla de resultados experimentales (convencion de latex-tables: booktabs,
% sin lineas verticales ni \hline). Rellenar SOLO desde data/stats_summary.json.
\begin{table}[htbp]
\centering
\caption{TODO: descripcion de la tabla de experimentos.}
\label{tab:experiment-results}
\begin{tabular}{lrr}
\toprule
Condicion & Metrica 1 & Metrica 2 \\
\midrule
TODO: condicion A & \todo{DATO PENDIENTE: metrica 1} & \todo{DATO PENDIENTE: metrica 2} \\
TODO: condicion B & \todo{DATO PENDIENTE: metrica 1} & \todo{DATO PENDIENTE: metrica 2} \\
\bottomrule
\end{tabular}
\end{table}

% TODO: analisis por hipotesis (H1, H2...) con IC 95% y tamano del efecto
% junto a cada contraste (ver .roo/skills/methods-reporting). El p-valor
% aislado no constituye evidencia suficiente.
